\section{Synthetic data generation approaches}
\label{sec:rq1}

Synthetic data generation methods differ in how they represent the process
that produces an observation. Statistical and learned models infer
relationships from examples, whereas simulation explicitly describes
scenes and their observation. Hybrid approaches combine these strategies.
This review uses these four groups to organise the literature, rather than
as mutually exclusive categories. They provide a basis for comparing how
generation methods control the resulting data, represent variation and
depend on real observations
\parencite{hernandezSyntheticDataGeneration2022,
nikolenkoSyntheticDataDeep2019,manivasagamLiDARsimRealisticLiDAR2020}.
Related \gls{augmentation} techniques transform or recombine existing
examples, so their contribution depends on both the original data and
whether the transformations preserve the required labels
\parencite{nikolenkoSyntheticDataDeep2019}.

\subsection{Statistical models}

Statistical generators produce samples from fitted distributions that
describe relationships between variables. For example,
\glspl{bayesian-network} represent conditional dependencies, allowing new
records to be sampled from the fitted model. The review by Hernandez et al.
and the benchmark by Hansen et al. discuss such methods for tabular data,
where observations consist of numerical or categorical attributes
\parencite{hernandezSyntheticDataGeneration2022,
hansenReimaginingSyntheticTabular2023}.
These studies establish a general approach to modelling data dependencies,
but do not demonstrate that tabular generators reproduce the spatial
structure of automotive LiDAR scans. Their relevance here is the underlying
principle that generated samples reflect the relationships represented by
the model. Increasing the number of samples cannot by itself recover
relationships that were omitted or poorly estimated, and uncertainty in
the fitted generator therefore remains relevant when synthetic datasets
are enlarged \parencite{breugelSyntheticDataReal2023}.

\subsection{Procedural and physical simulation}

Simulation provides a different form of control by specifying how scenes
are constructed and how a sensor observes them. Procedural rules define
scene content and variation, while the observation model converts that
content into measurements. This separation allows scene configurations to
be varied deliberately and annotations to be derived from known scene
properties. However, the generated observations remain constrained by the
available assets and the effects represented in the simulator
\parencite{nikolenkoSyntheticDataDeep2019,
manivasagamLiDARsimRealisticLiDAR2020}.
\Gls{domain-randomization} broadens the generated conditions by varying
selected parameters. Its usefulness consequently depends on which
parameters are varied and how their values are sampled, rather than on
variation alone \parencite{nikolenkoSyntheticDataDeep2019}.
For LiDAR perception, this distinction motivates examining scene diversity
and sensor modelling separately: changing the objects in a scene does not
necessarily correct how the simulated sensor samples their surfaces.

\subsection{Learned generators}

Learned generators model patterns in example data without requiring every
relationship to be specified explicitly. Approaches include
\glspl{vae}, which generate through a learned latent representation,
\glspl{gan}, which use adversarial training, and
\glspl{diffusion-model}, which learn a denoising process.
Hansen et al. compare representatives of these families in a tabular
setting, illustrating the range of available modelling approaches rather
than establishing a preferred method for LiDAR
\parencite{hansenReimaginingSyntheticTabular2023}.
Their flexibility also makes the training examples important: plausible
individual samples do not establish that a generator covers the full
distribution, and effects such as \gls{mode-collapse} can restrict the
variation represented \parencite{breugelSyntheticDataReal2023}.

LiDARGen provides a sensor-specific example of learned generation. It
represents scans as range and intensity images, learns through
\gls{score-matching}, and generates samples using iterative
\gls{langevin-dynamics}. It also supports conditional scan densification
\parencite{zyrianovLearningGenerateRealistic2022}.
This demonstrates how a learned model can represent scan structure, but
does not make scan generation equivalent to constructing an annotated
training dataset. Semantic labels and explicit control of scene objects
are additional requirements that cannot be inferred from plausible range
and intensity values alone. Learned generation must therefore be assessed
against the intended use of its outputs, including the annotations and
controls needed by the perception task.

\subsection{Hybrid approaches and implications for LiDAR}

Hybrid pipelines retain explicit scene or sensor structure while learning
components that are difficult to describe analytically. LiDARsim combines
assets reconstructed from real recordings with \gls{ray-casting} and a
learned \gls{ray-drop} model. The geometric model determines potential
surface returns, while the learned component models missing returns
\parencite{manivasagamLiDARsimRealisticLiDAR2020}.
This combination connects controllable scene construction with
data-dependent sensor modelling. It also shows why synthetic perception
training should not be interpreted as independence from real data:
recordings and annotations can still be required to reconstruct assets or
fit parts of the generator.

Together, these approaches reveal a recurring trade-off between explicit
control and dependence on observed examples. Statistical and learned
models reproduce relationships captured during fitting, simulation
exposes scene and sensor parameters, and hybrids combine explicit models
with learned components. These differences make control, coverage,
annotation availability and real-data requirements more informative
comparison criteria than dataset size alone. Applying this comparison to
automotive LiDAR requires considering how each approach constructs a
scene, samples visible surfaces and associates measurements with labels.
The next section examines these requirements in concrete LiDAR generation
pipelines.
